\documentclass[a4paper,11pt,twoside]{article}


\usepackage[utf8]{inputenc}
\usepackage[ngerman]{babel}
\usepackage{graphicx}
\usepackage{geometry}
\usepackage{setspace}
\usepackage{fancyhdr}
\usepackage{titlesec}
\usepackage{tocbibind}
\usepackage{csquotes}
\usepackage[backend=biber,style=authoryear]{biblatex}
\usepackage{hyperref}
\usepackage{helvet}
\renewcommand{\familydefault}{\sfdefault} % Serifenlose Schrift wie Arial
\usepackage[utf8]{inputenc}
\usepackage{caption}
\usepackage{tocloft} % Für Inhaltsverzeichnis-Anpassungen

% Seitenränder für doppelseitiges Layout
\geometry{top=2.5cm, bottom=2.5cm, left=2.5cm, right=2.5cm}

% Zeilenabstand
\onehalfspacing

% Kopf- und Fusszeilen mit fancyhdr
\fancypagestyle{standard}{
	\fancyhf{} % Leert die Kopf- und Fusszeilen
	% Kopfzeilen
	\fancyhead[LE]{\nouppercase{\leftmark}} % Kapitelname links auf geraden Seiten (kein ALL CAPS)
	\fancyhead[RO]{\nouppercase{\leftmark}} % Kapitelname rechts auf ungeraden Seiten (kein ALL CAPS)
	% Fusszeilen
	\fancyfoot[LE]{\thepage}                % Seitenzahl links auf geraden Seiten
	\fancyfoot[RO]{\thepage}                % Seitenzahl rechts auf ungeraden Seiten
}

\fancypagestyle{appendix}{
	\fancyhf{}
	% Kopfzeilen
	\fancyhead[LE]{Anhang}                  % "Anhang" links auf geraden Seiten
	\fancyhead[RO]{Anhang}                  % "Anhang" rechts auf ungeraden Seiten
	% Fuzeilen
	\fancyfoot[LE]{\thepage}                % Seitenzahl links auf geraden Seiten
	\fancyfoot[RO]{\thepage}                % Seitenzahl rechts auf ungeraden Seiten
}

\fancypagestyle{plain}{%
	\fancyhf{} % Leert die Kopf- und Fuzeilen
	\fancyhead[LE]{\nouppercase{\leftmark}} % Kapitelname links auf geraden Seiten
	\fancyhead[RO]{\nouppercase{\leftmark}} % Kapitelname rechts auf ungeraden Seiten
	\fancyfoot[LE]{\thepage} % Seitenzahl links auf geraden Seiten
	\fancyfoot[RO]{\thepage} % Seitenzahl rechts auf ungeraden Seiten
}

% Inhaltsverzeichnis anpassen
\renewcommand{\cfttoctitlefont}{\large\bfseries} % Titel Inhaltsverzeichnis in \large und fett
\renewcommand{\cftsecfont}{\normalfont} % Sektionen nicht fett
\renewcommand{\cftsubsecfont}{\normalfont} % Subsektionen nicht fett
\renewcommand{\cftsecpagefont}{\normalfont} % Seitenzahlen nicht fett
\renewcommand{\cftsubsecpagefont}{\normalfont} % Seitenzahlen von Subsektionen nicht fett

% Punkte für alle Ebenen im Inhaltsverzeichnis
\renewcommand{\cftsecleader}{\cftdotfill{\cftdotsep}} % Punkte für \section
\renewcommand{\cftsubsecleader}{\cftdotfill{\cftdotsep}} % Punkte für \subsection
\renewcommand{\cftsubsubsecleader}{\cftdotfill{\cftdotsep}} % Punkte für \subsubsection

% Einzug für jede Ebene anpassen
\setlength{\cftsecindent}{0em} % Keine Einrückung für \section
\setlength{\cftsubsecindent}{2em} % Einrückung für \subsection
\setlength{\cftsubsubsecindent}{4em} % Einrückung für \subsubsection

% Zeigt Subsections im Inhaltsverzeichnis an
\setcounter{tocdepth}{2}
% Optional: Nummeriert Subsections
\setcounter{secnumdepth}{2}

% Keine Kopf- und Fusszeilen vor der Einleitung
\pagestyle{empty} % "empty" deaktiviert die Kopf- und Fusszeilen

% Überschriftenformatierung
\titleformat{\section}{\normalfont\large\bfseries}{\thesection}{1em}{}
\titleformat{\subsection}{\normalfont\normalsize\bfseries}{\thesubsection}{1em}{}
\titleformat{\subsubsection}{\normalfont\normalsize\bfseries}{\thesubsubsection}{1em}{}

% Globale Einstellung für Tabellen-Captions
\captionsetup[table]{justification=raggedright, singlelinecheck=false}
\captionsetup{{labelfont=bf, small}, textfont=small}
% Tabellen-/Abbildungsbeschriftung kleiner und fettgedruckt, text nur small

% Hyperlinks
\hypersetup{
	colorlinks=true,
	linkcolor=black,
	citecolor=black,
	filecolor=black,
	urlcolor=blue,
	pdftitle={Titel der Arbeit},
	pdfauthor={Autor/in},
	pdfsubject={Art der Arbeit},
	pdfkeywords={Schlagworte}
}



% BibLaTeX-Bibliografie-Datei einbinden
\addbibresource{literatur.bib} % Pfad zur BibTeX-Datei

\begin{document}
	
	% Titelseite
	\begin{titlepage}
		\centering
		{\large \MakeUppercase{Zürcher Hochschule für Angewandte Wissenschaften}\\
			\MakeUppercase{Departement Life Sciences und Facility Management}\\
			\MakeUppercase{Institut für Umwelt und natürliche Ressourcen}\\[1.5cm]}
		
		% Titel der Arbeit
		{\Large \textbf{Titel der Arbeit}\\[1cm]}
		%{\large Untertitel\\[1.5cm]} falls Untertitel notwendig
		{\large{Art der Arbeit}\\[2cm]}
		
		% zentrierte Angaben
		\begin{center}
			\textbf{von}\\
			\textbf{Name, Vorname}\\
			Dein Studiengang\\
			Abgabedatum: TT.MM.JJJJ
		\end{center}
		
		%Fachkorrektur
		\vspace{3cm}
		\begin{tabbing}
			\textbf{Fachkorrektur:} \\
			(Akad. Titel) Name, Vorname (ausgeschrieben)\\
			Firma, Adresse, Ort\\
			\\%leere zeile
			(Akad. Titel) Name, Vorname (ausgeschrieben)\\
			Firma, Adresse, Ort\\
		\end{tabbing}
			
		% Logo der Hochschule
		\vfill
		\begin{figure}[t]
			\raggedleft
			\includegraphics[width=0.5\textwidth]{zhaw.jpg} % Logo-Datei
		\end{figure}
	\end{titlepage}
	
	
	% Impressum auf der Rückseite des Titelblatts
	\newpage
	\thispagestyle{empty}
	\vspace*{18cm}
	\noindent
	\textbf{Impressum}\\[1cm]
	\textbf{Schlagworte:} Schlagwort1, Schlagwort2, Schlagwort3\\[0.5cm]
	\textbf{Zitiervorschlag:}\\
	Autor/in (Jahr). Titel der Arbeit. Zürcher Hochschule für Angewandte Wissenschaften, Departement Life Sciences und Facility Management.
	
	
	% Zusammenfassung
	\newpage
	\section*{Zusammenfassung}
	blablabla
	
	% Abstract
	\newpage
	\section*{Abstract}
	blablabla
	
	% Abkürzungsverzeichnis
	\newpage
	\section*{Abkürzungsverzeichnis}
	\begin{tabbing}
		\hspace*{3cm} \= \kill
		Abb. \> Abbildung\\
		Tab. \> Tabelle\\
		etc.
	\end{tabbing}
	
	
	% Inhaltsverzeichnis
	\newpage
	\addtocontents{toc}{\protect\setcounter{tocdepth}{-1}} %kapitel Inhaltsverzeichnis ausschliessen
	\tableofcontents
	\addtocontents{toc}{\protect\setcounter{tocdepth}{3}}
	\thispagestyle{empty}
	
	
	
	% Einleitung
	\newpage
	\pagestyle{standard} % Aktiviert Kopf- und Fußzeilen ab dieser Seite
	\section{Einleitung}
	blablabla
	\subsection{Stand der Forschung}
	blablabla
	\subsubsection{Unterunterkapitel}
	blablabla \parencite{gislerProjectWorkResearch2023}
	\\
	\textcite{wust-galleyModelbasedCarbonInventory2020}
	\\
	\
	
	% Weitere Kapitel
	\newpage
	\section{Theorie oder Literaturübersicht}
	
	\begin{figure}[b] %[ht] nach \begin{figure}, um die Abbildung "hier" (h) oder oben (t) zu platzieren. Andere Optionen sind b (bottom) oder p (auf einer separaten Seite)
		\raggedright %\centering, \raggedright oder \raggedleft für Position
		\includegraphics[width=0.5\textwidth]{zhawGR_logo.png} % Ersetzt "beispielbild" durch den Namen der Bilddatei (z. B. beispielbild.png), muss im gleichen Verzeichnis sein
		\caption{Beispielabbildung}
		\label{fig:beispiel}
	\end{figure}
	Abbildung \ref{fig:beispiel} zeigt ein Beispiel.
	
	
	\newpage
	\section{Material und Methoden}
	\begin{table}[ht]
		\raggedright
		\caption{Beispieltabelle}
		\begin{tabular}{|c|c|c|}
			\hline
			Spalte 1 & Spalte 2 & Spalte 3 \\ \hline
			Daten 1  & Daten 2  & Daten 3  \\ \hline
			Daten 4  & Daten 5  & Daten 6  \\ \hline
		\end{tabular}
		\label{tab:beispiel}
	\end{table} 
	\noindent %kein Einzug
	blablabla Tabelle \ref{tab:beispiel} zeigt ein Beispiel.
	\\%neue Zeile
	\par\noindent neuer Absatz ohne Einzug
	\begin{table}[ht]
		\centering
		\begin{tabular}{|c|c|c|}
			\hline
			Spalte 1 & Spalte 2 & Spalte 3 \\ \hline
			Daten 1  & Daten 2  & Daten 3  \\ \hline
			Daten 4  & Daten 5  & Daten 6  \\ \hline
		\end{tabular}
		\caption{andere Tabelle}
		\label{tab:andere}
	\end{table}
	\par\noindent  Tabelle \ref{tab:andere} zeigt ein Beispiel.
	
	\newpage
	\section{Ergebnisse}
	% Inhalt
	
	\newpage
	\section{Diskussion}
	blablabla
	% Inhalt
	
	
	% Literaturverzeichnis
	\newpage
	\pagestyle{plain} % Setzt den neuen Seitenstil
	\section*{Literaturverzeichnis}
	\printbibliography[heading=none]
	\markboth{Literaturverzeichnis}{Literaturverzeichnis} % Setzt den Titel in die Kopfzeile
	\addcontentsline{toc}{section}{Literaturverzeichnis}
	
	
	% Abbildungsverzeichnis (falls benötigt)
	\newpage
	\pagestyle{plain} % Setzt den neuen Seitenstil
	\listoffigures
	\markboth{Abbildungsverzeichnis}{Abbildungsverzeichnis} % Kapitelname in der Kopfzeile
	
	
	% Tabellenverzeichnis (falls benötigt)
	\newpage
	\pagestyle{plain} % Setzt den neuen Seitenstil
	\listoftables
	\markboth{Tabellenverzeichnis}{Tabellenverzeichnis} % Kapitelname in der Kopfzeile
	
	
	% Anhang (falls benötigt)
	\newpage
	\appendix
	\cleardoublepage %falls anhang mit ungerader seite beginnen muss
	\pagenumbering{arabic} % Seitenzahlen im Anhang beginnen bei 1
	\pagestyle{appendix}   % Seitenstil für den Anhang (gerade/ungerade)
	\section*{Anhang}
	\addcontentsline{toc}{section}{Anhang}
	\renewcommand{\thesubsection}{Anhang \Alph{subsection}}
	\markboth{Anhang}{Anhang}
	asdf
	\newpage
	\noindent asdf
	
	
	
\end{document}
